% !TeX encoding = UTF-8
\SetChapterInfo{subtitle={从时间平均到扩散生成元与抽象位势},keywords={Markov; 遍历; 熵; Brownian; Feller; 跳跃; 位势},prereq={平均遍历; 正算子; 强连续半群; 向量格; 测度表示},goals={区分平均与逐点收敛; 解释局部与跳跃生成元; 控制位势定义域}}
\chapter[遍历理论和扩散理论]{遍历理论\\和扩散理论}\label{ch:yosida-13}
\ChapterSource{按 \cite{yosida-zh,yosida-en} 第十三章九节组织，2026 年 10 月 27 日补写。直接查看英文印刷 379--418 页全部扫描图，第 418 页上部仍为本章结尾、下部开始第十四章，故修正本章英文范围。中译本范围为 397--436 页，选页与证明边界详见本章 COVERAGE。}

遍历理论研究长时间平均，扩散理论研究短时间变化；半群把两端联系起来。不变测度产生各 $L^p$ 空间上的收缩，极大不等式使平均收敛提升为几乎处处收敛；局部矩与跳跃测度决定生成式，无折扣预解式极限则产生位势与 Poisson 方程。

遍历部分先用概率空间与固定一步算子 $T=T_1$；连续时间结论另列假设。记 $A_n=n^{-1}(I+T+\cdots+T^{n-1})$，它与生成元 $A$ 的记号按上下文区分。扩散部分采用 $b\cdot\nabla+\tfrac12Q:D^2$，热核模型的生成元为 $c\Delta$。抽象位势统一取 $V=-A^{-1}$，以 $\lambda\downarrow0$ 的 Abel 极限定义。

\Needspace{10\baselineskip}
\section{具有不变测度的 Markov 过程}\label{sec:yosida-13-01}
\noindent\textit{原书起始页：中译本 397；英文版 379。}

\begin{definition}[转移核]
可测空间 $(S,\Sigma)$ 上的概率转移核 $P$ 满足：$P(x,\cdot)$ 为概率测度，$P(\cdot,B)$ 可测。时间齐次转移族满足 $P_0(x,\cdot)=\delta_x$ 以及 Chapman--Kolmogorov 方程
\[
 P_{t+s}(x,B)=\int_S P_s(y,B)P_t(x,\dd y).
\]
\end{definition}

测度 $m$ 不变，指 $\int P_t(x,B)\,m(\dd x)=m(B)$。有限非零不变测度可归一化，本章遍历部分约定 $m(S)=1$。记
\[
 (T_tf)(x)=\int_S f(y)P_t(x,\dd y).
\]

本章主要研究转移核与其算子。转移族的存在不自动等于已构造出具有连续样本路径的随机过程；样本路径构造所需的概率空间与正则性应另外说明。

\begin{proposition}[Markov 收缩]\label{prop:y13-markov-contraction}
若 $m$ 不变，则 $T_t$ 在每个 $L^p(m)$（$1\le p\le\infty$）上诱导正收缩，满足 $T_t1=1$、$\int T_tf\,\dd m=\int f\,\dd m$ 以及半群性质。
\end{proposition}
\begin{proof}
对 $m$ 零集 $N$，$\int P_t(x,N)\,\dd m=0$，故 $P_t(x,N)=0$ 几乎处处；算子因而与代表元无关。对有界函数由核积分得到正性、单位与保积分；单调收敛延伸到非负可积函数。Jensen 不等式给出 $|T_tf|^p\le T_t|f|^p$，积分并用不变性即得有限 $p$ 的估计。$p=\infty$ 用概率加权平均估计。截断及完备性将算子延拓到各 $L^p$；Fubini 与 Chapman--Kolmogorov 方程给出半群性质。
\end{proof}

\begin{theorem}[有限不变测度下的平均遍历]\label{thm:y13-markov-mean}
固定一步转移 $T=T_1$，置 $A_n=n^{-1}\sum_{k=0}^{n-1}T^k$。对每个 $f\in L^p(m)$、$1\le p<\infty$，$A_nf$ 在 $L^p$ 中收敛到 $Qf$，满足
\[
 TQf=Qf,\quad Q^2=Q,\quad Q1=1,\quad
 \int Qf\,\dd m=\int f\,\dd m.
\]
$Q$ 为正收缩投影。若 $T_t$ 在 $L^p$ 上强连续，相同结论适用于 $t^{-1}\int_0^tT_sf\,\dd s$，其极限投到全半群的共同不动空间。
\end{theorem}
\begin{proof}
在 $L^2$ 上应用第八章的 Hilbert 空间平均遍历定理。对 $L^1$，有界函数在其中稠密，且 $\|g\|_1\le\|g\|_2$。先对有界函数得到 $L^1$ 收敛，再用 $\sup_n\|A_n\|_1\le1$ 延拓到全部 $L^1$。对有界 $g$，$A_ng$ 一致有界，故其 $L^1$ 收敛通过 $|h|^p\le C^{p-1}|h|$ 提升到任意有限 $p$；再次使用有界函数的稠密性完成证明。

$TA_n-A_n=(T^n-I)/n$，所以极限不动；不动元素被每个平均保持，故 $Q^2=Q$。其余性质由极限得到。连续时间在 $L^2$ 上使用连续平均遍历，再作同样稠密延拓；端段积分估计给出对所有 $T_s$ 的不变性。
\end{proof}

\begin{example}[确定性过程与有限状态链]
若可测变换 $\tau$ 保持 $m$，则 $P(x,B)=\mathbf1_B(\tau x)$，$Tf=f\circ\tau$，是 Markov 核的特殊情形。有限状态时 $P=(p_{ij})$ 为行随机矩阵，不变概率是行向量 $\pi P=\pi$；算子在函数列向量上左乘 $P$，不要把两个方向混用。
\end{example}

\begin{exercise}[周期链]\label{ex:y13-periodic-chain}
对 $P=\left(\begin{smallmatrix}0&1\\1&0\end{smallmatrix}\right)$ 求不变概率和 $A_nf$ 的极限，并说明 $P^nf$ 未必收敛。
\end{exercise}
\begin{solution}
不变概率为 $(1/2,1/2)$。常数方向特征值为 $1$，$(1,-1)$ 方向特征值为 $-1$，故 $A_n(f_1,f_2)\to((f_1+f_2)/2,(f_1+f_2)/2)$。在第二个方向，$P^n$ 交替变号。平均收敛与各时刻的分布收敛是不同命题。
\end{solution}

\begin{exercise}[不变性不可省略]\label{ex:y13-noninvariant}
在两点空间取均匀概率，令两行转移概率均为 $(1,0)$。证明该算子保持常数，却不是 $L^1$ 收缩。
\end{exercise}
\begin{solution}
$T(a,b)=(a,a)$，故 $T1=1$。取 $f=(1,0)$，则 $\|f\|_1=1/2$ 而 $\|Tf\|_1=1$。均匀概率不满足 $\pi P=\pi$，不能从 Markov 性单独推得相对于给定测度的 $L^1$ 收缩。
\end{solution}

\Needspace{10\baselineskip}
\section{个别遍历定理及其应用}\label{sec:yosida-13-02}
\noindent\textit{原书起始页：中译本 401；英文版 383。}

本节 $T$ 为上一节的 Markov 收缩，$m(S)=1$，$A_n$ 从 $k=0$ 起求和。将原书从 $k=1$ 起的记法改到此处约定，不改变下述极限。证明采用有限极大不等式；它把稠密子空间上的逐点收敛延拓到 $L^1$，避免误把稠密集视为非第一纲集。

\begin{lemma}[Hopf 有限极大不等式]\label{lem:y13-hopf-maximal}
对实 $f\in L^1$，令 $S_kf=\sum_{j=0}^{k-1}T^jf$，$E_N=\{\max_{1\le k\le N}S_kf>0\}$。则 $\int_{E_N}f\,\dd m\ge0$。因而对任意 $f\in L^1$ 和 $\alpha>0$，
\begin{equation}\label{eq:y13-maximal-weak}
 m\{\sup_{n\ge1}|A_nf|>\alpha\}\le\|f\|_1/\alpha.
\end{equation}
\end{lemma}
\Needspace{4\baselineskip}
\begin{proof}
设 $g_N=\max(0,S_1f,\ldots,S_Nf)$、$g_0=0$。由于 $T$ 正，$S_kf=f+TS_{k-1}f\le f+Tg_{N-1}$。在 $E_N$ 上取最大值得 $g_N\le f+Tg_{N-1}$，故
\[
 \int_{E_N}f\ge\int g_N-\int_{E_N}Tg_{N-1}
 \ge\int g_N-\int g_{N-1}\ge0.
\]
对 $|f|-\alpha$ 应用此式；$T1=1$ 使正部分和事件恰为 $\max_{k\le N}A_k|f|>\alpha$。于是 $\alpha m(E_N)\le\|f\|_1$。用 $|A_nf|\le A_n|f|$ 并令 $N\to\infty$，得到所需估计。
\end{proof}

\begin{theorem}[Markov 算子的个别遍历定理]\label{thm:y13-individual}
对每个 $f\in L^1(m)$，$A_nf$ 几乎处处收敛到有限函数 $Qf$，且在 $L^1$ 中收敛。对 $f\in L^p$、$1\le p<\infty$，极限也与上一节的 $L^p$ 极限相同。
\end{theorem}
\begin{proof}
在 Hilbert 空间 $L^2$ 中，收缩性给出 $\ker(I-T)=\ker(I-T^*)$，因此
\[
 L^2=\ker(I-T)\oplus\overline{(I-T)L^2}.
\]
为说明第一个等式，若 $Th=h$，则 $\|T^*h\|\le\|h\|$ 且 $\langle T^*h,h\rangle=\|h\|^2$，Cauchy--Schwarz 的等号情形给出 $T^*h=h$；反向同理。由于有界函数在 $L^2$ 中稠密，
\[
 \mathcal D=\ker(I-T)|_{L^2}+(I-T)L^\infty
\]
在 $L^2$ 及 $L^1$ 中稠密。对不动函数，平均本来不变；对 $g=(I-T)h$、$h\in L^\infty$，有 $A_ng=(h-T^nh)/n\to0$ 几乎处处，甚至本质一致。因此 $\mathcal D$ 中每个元素都有有限逐点极限。

先取实 $f\in L^1$。极大估计在 $\alpha\to\infty$ 时说明 $\sup_n|A_nf|<\infty$ 几乎处处。其上、下极限之差记为 $\omega(f)$。任取 $g\in\mathcal D$，由于 $A_ng$ 收敛，
\[
 \omega(f)\le2\sup_n A_n|f-g|,
 \qquad m\{\omega(f)>2\alpha\}\le\|f-g\|_1/\alpha.
\]
令 $g\to f$，得到 $\omega(f)=0$ 几乎处处。复函数分别处理实虚部。上一节已有 $L^1$ 极限；抽取几乎处处收敛的子序列后，极限的唯一性识别该函数为 $Qf$。有限 $p$ 情形亦然。
\end{proof}

这与第十二章的 Banach 收敛思想一致：关键是先证明控制收敛集闭性的估计。仅在一个稠密子空间验证收敛并不足够。对连续时间，本节结论先用于一步平均 $g=\int_0^1T_sf\,\dd s$ 得整数时刻的积分平均；一般几乎处处连续时间版本还需共同可测代表与余项控制，此处不据此省略这些条件。

\Needspace{7\baselineskip}
\paragraph{原书的比值推广。}
Chacon--Ornstein 比值遍历定理对一般正 $L^1$ 收缩 $T$ 断言：若 $f,p\in L^1$、$p\ge0$，则
\[
 \frac{\sum_{k=0}^{n-1}T^kf}{\sum_{k=0}^{n-1}T^kp}
\]
在 $\{\sum_{k\ge0}T^kp>0\}$ 上几乎处处有有限极限；分母最终为正后才定义比值。本节完整证明的是上述保概率 Markov 情形；一般比值定理及原书的消去引理作为明确标注的外部推广保留，不将极大不等式的证明说成已证明整个比值理论。

\begin{exercise}[上同调差的平均]\label{ex:y13-coboundary}
若 $g=(I-T)h$ 且 $h\in L^\infty$，求 $A_ng$ 的显式误差界。说明同一公式在 $h\in L^1$ 时只立即给出哪种收敛。
\end{exercise}
\begin{solution}
$A_ng=(h-T^nh)/n$，故 $\|A_ng\|_\infty\le2\|h\|_\infty/n$。若只有 $h\in L^1$，同样得到 $\|A_ng\|_1\le2\|h\|_1/n$，但该范数估计本身不推出全序列几乎处处收敛；本节定理另给出后者。
\end{solution}

\begin{exercise}[保持变换的时间平均]\label{ex:y13-birkhoff}
若 $\tau$ 保持概率测度且 $f\in L^1$，把本节定理写成轨道函数公式，并说明例外集对一个给定 $f$ 的含义。
\end{exercise}
\begin{solution}
取 $Tf=f\circ\tau$，则 $n^{-1}\sum_{k=0}^{n-1}f(\tau^kx)\to Qf(x)$ 几乎处处，并且积分保持。可为这个 $f$ 选一个零测例外集；可数个给定函数可以统一例外集。不能在任意不可数函数族上直接取零测集的并，声称得到所有可积函数同时逐点收敛。
\end{solution}

\Needspace{10\baselineskip}
\section{遍历假设和 \texorpdfstring{$H$}{H}-定理}\label{sec:yosida-13-03}
\noindent\textit{原书起始页：中译本 407；英文版 389。}

\begin{definition}[遍历与混合]
保持概率 $m$ 的 Markov 算子称为遍历的，若每个 $f\in L^1$ 的时间平均均趋于常数 $\int f\,\dd m$。集合 $B$ 称为不变集，若 $T\mathbf1_B=\mathbf1_B$ 几乎处处。不变集构成的完备化 $\sigma$ 代数记为 $\mathcal I$。若对每对 $f,g\in L^2$ 都有
\[
 \langle T^nf,g\rangle\longrightarrow
 \left(\int f\,\dd m\right)\overline{\left(\int g\,\dd m\right)},
\]
则称一步算子混合；连续半群用 $t\to\infty$ 作相应定义。
\end{definition}

\begin{theorem}[遍历性的等价表述]\label{thm:y13-ergodic-equivalence}
下列条件等价：遍历性；$\ker(I-T)|_{L^2}$ 只有常数；每个不变集的概率为 $0$ 或 $1$；对任意可测集 $B,C$，
\[
 \frac1n\sum_{k=0}^{n-1}\int_B P^k(x,C)\,m(\dd x)
 \longrightarrow m(B)m(C).
\]
一般情形下，平均投影是条件期望 $Qf=\mathbb E_m(f\mid\mathcal I)$。
\end{theorem}
\begin{proof}
前两条件由平均投影的值域及 $L^2$ 在 $L^1$ 中稠密得到。若所有不动函数为常数，不变集的示性函数只能为零或一。反向先考虑实不动函数 $f\in L^2$。Jensen 及不变性给出
\[
 T(f^2)-(Tf)^2\ge0,\qquad \int[T(f^2)-f^2],\dd m=0.
\]
故对几乎每个 $x$，$P(x,\cdot)$ 集中于 $f(y)=f(x)$ 的集合。于是每个水平集 $\{f>a\}$ 都不变。若不变集概率只能为零或一，$f$ 必为常数。水平集判别也说明一般不动空间恰为 $L^2(\mathcal I)$；不变集在核下几乎必然保持其内外，故确实对补及可数并封闭。

第一节证明的 $\ker(I-T)=\ker(I-T^*)$ 表明不变示性函数也被 $T^*$ 保持。因此 $\int_B A_nf=\int_B f$ 对 $B\in\mathcal I$ 成立，取极限得到条件期望的刻画，先在 $L^2$ 后在 $L^1$ 成立。最后取 $f=\mathbf1_C$、$g=\mathbf1_B$ 得集合相关式；反向由简单函数稠密性及收缩性，将该式延伸到任意 $L^2$ 函数，识别平均投影只有常数。
\end{proof}

混合蕴含遍历，因为普通极限必有相同的 Ces\`aro 极限；反向不成立。还须区分固定一步 $T_1$ 的遍历与全连续半群的遍历：二者的不动空间可能不同，离散采样可产生共振。

\begin{example}[无理旋转遍历但不混合]
圆周上 $\tau(x)=x+\alpha\pmod1$，$\alpha$ 无理。Fourier 模 $e_k(x)=e^{2\pi ikx}$ 满足 $Te_k=e^{2\pi ik\alpha}e_k$，不动函数只有常数，故遍历。但 $|\langle T^ne_1,e_1\rangle|=1$，所以不混合。更一般地，遍历 Markov 算子的单位模本征函数满足 $|f|$ 为常数：$T|f|\ge|Tf|=|f|$，积分相等迫使等号，再用不动空间的刻画。这正是原书角变量的来源。
\end{example}

\Needspace{14\baselineskip}
\begin{theorem}[$H$-定理：熵的单调性]\label{thm:y13-entropy}
若 $f\ge0$、$\int f\,\dd m=1$ 且 $\int f\log^+f\,\dd m<\infty$，定义
\[
 \mathcal H(f)=-\int f\log f\,\dd m,\qquad 0\log0=0.
\]
则 $\mathcal H(T_tf)\ge\mathcal H(f)$，且对 $t_2\ge t_1$ 有 $\mathcal H(T_{t_2}f)\ge\mathcal H(T_{t_1}f)$。相反号的相对熵 $\int f\log f$ 单调不增。
\end{theorem}
\begin{proof}
$\Phi(r)=r\log r$ 凸且下界为 $-e^{-1}$，核的 Jensen 不等式给出 $\Phi(T_tf)\le T_t\Phi(f)$。通过截断凸函数或积分形式的 Jensen 可处理无界 $f$；右侧可积，且不变性给出 $\int\Phi(T_tf)\le\int\Phi(f)$。再用半群性质得到随时间的单调性。该结论只利用保积分与 Markov 性，不依赖遍历性，也不自动给出熵趋于最大值。
\end{proof}

这里证明的是 $T_t$ 对非负函数的 Jensen 不等式。若把 $f$ 解释为给定过程的分布密度，其前向演化通常由核的对偶决定，须在相应对偶空间中另行核对上述性质。

\begin{exercise}[采样共振]\label{ex:y13-sampling}
圆周上的连续平移 $T_tf(x)=f(x+t)$ 的共同不动函数有哪些？比较 $T_1$ 的不动空间。
\end{exercise}
\begin{solution}
所有 $t\ge0$ 的平移共同不变时，Fourier 系数中只有零频可以非零，故为常数。但 $T_1=I$，其不动空间为整个 $L^2$。连续时间平均投到常数，整数时间采样平均却保持原函数。
\end{solution}

\begin{exercise}[熵不保证混合]\label{ex:y13-entropy-rotation}
若 $T$ 来自保测变换，证明每个具有有限熵的概率密度满足 $\mathcal H(Tf)=\mathcal H(f)$。结合无理旋转说明熵单调性的限度。
\end{exercise}
\begin{solution}
因为 $Tf=f\circ\tau$，$\Phi(Tf)=\Phi(f)\circ\tau$，保测性使两者积分相等。无理旋转遍历却不混合，熵在每一步都保持；熵单调并不意味着严格增长、分布收敛或混合。
\end{solution}

\Needspace{10\baselineskip}
\section{带局部紧相空间的 Markov 过程的遍历分解}\label{sec:yosida-13-04}
\noindent\textit{原书起始页：中译本 411；英文版 393。}

以下先固定一步核 $P$。设 $S$ 为局部紧、第二可数 Hausdorff 空间，$C_0(S)$ 表示无穷远消失的连续函数。假定 $T C_0(S)\subset C_0(S)$；这是一项 Feller 性质，不由核的可测性自动推出。$C_c(S)$ 则表示紧支集连续函数，勿与旧文献中相近的记号混用。

\begin{theorem}[不变测度与不逃逸的平均质量]\label{thm:y13-invariant-existence}
存在非零有限不变正测度，当且仅当某个点 $x$ 与非负函数 $g\in C_0(S)$ 满足
\[
 \limsup_n A_ng(x)>0.
\]
若对每个 $x$ 和每个 $f\in C_0(S)$ 都有 $A_nf(x)\to0$，则不存在这样的不变测度。
\end{theorem}
\begin{proof}
令 $\eta_n=n^{-1}\sum_{k=0}^{n-1}P^k(x,\cdot)$。这些概率测度作为 $C_0(S)$ 上的泛函，其范数为一。$C_0(S)$ 可分，用可数稠密集及对角抽取得到某子序列在所有 $C_0$ 函数上收敛到正泛函 $L$；相应正则测度 $\eta$ 满足 $\eta(S)\le1$。选择子序列时保留 $\eta_n(g)$ 的正下界，故 $\eta\ne0$。由于 $Tf\in C_0$ 且
\[
 \eta_n(Tf)-\eta_n(f)=\frac{T^nf(x)-f(x)}n\longrightarrow0,
\]
$\eta$ 不变。这里使用 Riesz--Markov 表示及测度由 $C_0$ 积分确定的性质。

反之，若 $\eta$ 为非零有限不变测度，选非负 $g\in C_0$ 使 $\eta(g)>0$。第二节应用于归一化的 $\eta$，给出 $A_ng\to Qg$ 几乎处处且 $\int Qg\,\dd\eta=\eta(g)>0$，所以至少有一点的极限为正。最后的断言也可直接由控制收敛与不变性推出。
\end{proof}

非紧空间中上述子序列只有弱星即模糊收敛，极限可能丢失总质量；故必须先证明它非零，再归一化。紧空间时常数 $1\in C(S)$，极限质量仍为一，这就是 Krylov--Bogoliubov 平均构造的离散版本。

\begin{theorem}[紧可度量 Feller 情形的遍历分解]\label{thm:y13-ergodic-decomposition}
设 $S$ 紧可度量、$TC(S)\subset C(S)$，$m$ 为不变概率。存在可测概率核 $x\mapsto m_x$，使对 $m$ 几乎每个 $x$，$m_x$ 为遍历不变概率，并有
\[
 Qf(x)=\int f(y)m_x(\dd y)\quad(f\in C(S)),\qquad
 m(B)=\int m_x(B)m(\dd x).
\]
第一式可选共同的满测集，对所有连续函数同时成立。
\end{theorem}
\begin{proof}
取包含常数的可数实函数族，其线性包在 $C(S,\mathbb R)$ 中稠密。第二节给出一个共同满测集 $N$，在其中所有这些函数的平均收敛。由于 $\|A_n\|_\infty\le1$，对每个 $x\in N$，极限延拓为 $C(S)$ 上的正单位泛函 $L_x$。Riesz--Markov 给出概率 $m_x$。同样的端项估计说明 $L_x(Tf)=L_x(f)$，故 $m_x$ 不变。对连续函数，$x\mapsto m_x(f)$ 可测；单调类论证将其延伸为 Borel 概率核。积分平均保持性给出重心公式，也说明 $m_x(N)=1$ 对几乎每个 $x$ 成立。零测补集上的核可任取一点质量。

还须证明分量遍历，不能只从重心公式推断。记可数稠密族为 $f_j$，$q_j=Qf_j$，取上述极限代表。由第三节 $Q=\mathbb E(\cdot\mid\mathcal I)$，故
\[
 Qq_j=q_j,\qquad Q(q_j^2)=q_j^2.
\]
重心核在所有有界 Borel 函数上代表条件期望，因此
\[
 \int(q_j(y)-q_j(x))^2m_x(\dd y)=0
\]
对几乎每个 $x$、所有 $j$ 成立。于是 $m_x$ 集中在 $N$ 中使全部 $q_j(y)=q_j(x)$ 的共同水平集上。对其中的 $y$，$A_nf_j(y)\to q_j(x)=m_x(f_j)$。由于一致有界，控制收敛给出相对于 $m_x$ 的 $L^2$ 收敛到常数。连续函数在 $L^2(m_x)$ 中稠密，平均收缩使全部 $L^2(m_x)$ 函数的平均都趋于其积分，故 $m_x$ 遍历。
\end{proof}

一般局部紧空间中的完整“耗散部分加遍历部分”的逐点划分还涉及模糊极限的质量及共同例外集。上面完整证明了不变测度的存在判别与紧可度量情形的概率分解；原书更一般的全局划分及连续时间共同分解不在此处无条件推广。

\begin{exercise}[向无穷远逃逸]\label{ex:y13-escape}
在 $S=[0,\infty)$ 上取 $\tau(x)=x+1$。证明它保持 $C_0(S)$，且对每个 $f\in C_0(S)$、每个 $x$ 有 $A_nf(x)\to0$，故无不变概率。
\end{exercise}
\begin{solution}
$Tf(x)=f(x+1)$ 连续且无穷远消失。固定 $x$，$f(x+k)\to0$，其 Ces\`aro 平均也趋零。定理\ref{thm:y13-invariant-existence}排除非零有限不变测度。这说明非紧空间中的平均概率质量可能全部逃逸。
\end{solution}

\begin{exercise}[两个封闭类]\label{ex:y13-closed-classes}
四点链的转移矩阵为两个 $\left(\begin{smallmatrix}0&1\\1&0\end{smallmatrix}\right)$ 的直和。求全部不变概率、遍历不变概率及分解。
\end{exercise}
\begin{solution}
不变概率为 $(a/2,a/2,(1-a)/2,(1-a)/2)$，$0\le a\le1$。两个封闭类上的均匀概率 $m_1,m_2$ 分别遍历，全部不变概率为 $am_1+(1-a)m_2$。条件分量 $m_x$ 在第一个类等于 $m_1$，第二个类等于 $m_2$；周期性不妨碍各类的平均遍历性。
\end{solution}

\Needspace{10\baselineskip}
\section{齐次 Riemann 空间上的 Brown 运动}\label{sec:yosida-13-05}
\noindent\textit{原书起始页：中译本 416；英文版 398。}

Riemann 空间在本章理解为 Riemann 流形。转移核与二阶微分算子的联系，来自短时间内位移的一阶、二阶矩，而不是单靠“随机运动”一词。

\begin{definition}[局部性条件与空间齐次]
若对每个 $\varepsilon>0$、固定 $x$ 有
\[
 t^{-1}P_t(x,\{y:d(x,y)>\varepsilon\})\longrightarrow0,
\]
称满足 Lindeberg 型小时间条件；一致版本要求在所讨论的 $x$ 集合上一致。若传递等距群 $G$ 满足 $P_t(gx,gB)=P_t(x,B)$，称转移在空间上齐次。该条件使算子与群作用交换。
\end{definition}

\begin{proposition}[从局部矩得到生成元]\label{prop:y13-local-moments}
在点 $x$ 附近取坐标 $z^1,\ldots,z^d$，令 $\Delta z^i=z^i(y)-z^i(x)$。设上述小时间条件成立，截断一阶矩及二阶矩满足
\[
 t^{-1}\int_{d(x,y)<\varepsilon}\Delta z^i\,P_t(x,\dd y)\to b^i(x),\qquad
 t^{-1}\int_{d(x,y)<\varepsilon}\Delta z^i\Delta z^j\,P_t(x,\dd y)\to q^{ij}(x),
\]
极限不依赖充分小的 $\varepsilon$，且局部二阶矩除以 $t$ 一致有界。若 $f$ 有界、在 $x$ 附近为 $C^2$，则点态生成元为
\[
 Af(x)=\sum_i b^i(x)\partial_i f(x)+\tfrac12\sum_{i,j}q^{ij}(x)\partial_i\partial_j f(x).
\]
矩阵 $q$ 半正定。要把点态极限识别为某 Banach 空间生成元，还需相应范数收敛与定义域条件。
\end{proposition}
\begin{proof}
将积分分成小球内外。外部的差商绝对值不超过 $2\|f\|_\infty P_t(x,B_\varepsilon^c)/t$，趋零。内部作二阶 Taylor 展开，余项为 $o(|\Delta z|^2)$。先取 $t\downarrow0$，再令 $\varepsilon\downarrow0$，二阶矩的有界性使余项消失，主项即上式。对实向量 $\xi$，截断积分 $t^{-1}\int(\sum_i\xi_i\Delta z^i)^2\,\dd P_t\ge0$，取极限得到半正定性。原书把二阶系数直接记作 $b^{ij}$；此处以 $q/2$ 保留协方差的通常约定。
\end{proof}

\begin{example}[欧氏 Brownian 半群的完整模型]\label{exm:y13-gaussian-semigroup}
在 $\mathbb R^d$ 上取
\[
 p_t(z)=(4\pi ct)^{-d/2}e^{-|z|^2/(4ct)},\qquad T_tf=p_t*f,\quad c>0.
\]
Gaussian 卷积给出 $p_t*p_s=p_{t+s}$，总质量为一，故为正收缩半群。在 $C_0(\mathbb R^d)$ 上，近似单位估计给出强连续性；在 $C_c^\infty$ 上 Taylor 展开与 Gaussian 尾部估计给出 $Af=c\Delta f$，且差商在一致范数中收敛。位移均值为零、协方差为 $2ctI$；尾概率以指数速度衰减，满足 Lindeberg 条件。此为生成元 $c\Delta$ 的约定；常见概率论约定 $\tfrac12\Delta$ 对应 $c=1/2$。
\end{example}

\begin{proposition}[各向同性下的唯一微分形式]\label{prop:y13-isotropic}
若球面 $S^d$ 上的保守 Markov 生成元具有局部二阶形式，且与所有等距变换交换，则它为 $c\Delta_{S^d}$，$c\ge0$。在欧氏空间中相同结论适用于与全部欧氏等距变换交换的局部扩散生成元。
\end{proposition}
\Needspace{4\baselineskip}
\begin{proof}
在一点的法坐标中，稳定该点的等距群在切空间上包含全正交群。其不变向量只能为零，故漂移为零；其不变对称二阶型只能是单位型的数倍。正最大值性质使该倍数非负。传递性使系数在所有点相同，保常数性排除零阶项，故得到 Laplace--Beltrami 算子。
\end{proof}

一般齐次流形上的不变二阶型未必只有一个，因此“齐次”本身不能推出唯一的 $c\Delta$。一般流形热核的存在、保总质量及连续路径构造属于额外几何/PDE/概率理论，本节不把局部矩计算当成这些存在性定理的证明。紧连通流形上标准热半群保持归一化体积；其长期行为可再结合 Laplace 谱论分析。

\begin{exercise}[生成元的归一化]\label{ex:y13-brownian-scale}
一维核 $(2\pi t)^{-1/2}e^{-z^2/(2t)}$ 的生成元是什么？比较其二阶矩与 $p_t$ 在 $c=1$ 时的二阶矩。
\end{exercise}
\begin{solution}
该核方差为 $t$，故生成元为 $\tfrac12\,\dd^2/\dd x^2$。$c=1$ 的核方差为 $2t$，生成元为 $\dd^2/\dd x^2$。Taylor 展开的二分之一正是两个约定相差的来源。
\end{solution}

\begin{exercise}[平移不等于各向同性]\label{ex:y13-anisotropic}
对常向量 $b$ 和对称正定矩阵 $Q$，说明 $Af=b\cdot\nabla f+\tfrac12 Q:D^2f$ 的 Gaussian 半群为何平移齐次；何时还与全部正交变换交换？
\end{exercise}
\begin{solution}
转移为从 $x$ 出发、均值 $x+tb$、协方差 $tQ$ 的 Gaussian 分布，只依赖位移，故平移齐次。与所有正交变换交换要求 $Rb=b$、$RQR^{\mathsf T}=Q$ 对所有 $R$ 成立，所以 $b=0$、$Q=\sigma^2I$。平移齐次允许漂移与各向异性。
\end{solution}

\Needspace{10\baselineskip}
\section{W. Feller 的广义 Laplace 算子}\label{sec:yosida-13-06}
\noindent\textit{原书起始页：中译本 422；英文版 403。}

本节在连续函数空间上理解生成元，差商的一致范数收敛保证可逐点计算；常数不在空间中时，保守性由转移核表述。

\begin{proposition}[最大值性质与局部性]\label{prop:y13-local-maximum}
保常数的 Markov 生成元在其定义域上满足正最大值原理：若 $f$ 在 $x_0$ 达到非负全局最大值，则 $Af(x_0)\le0$。若另有上一节的 Lindeberg 条件，且 $f$ 在 $x_0$ 的邻域恒为零，则 $Af(x_0)=0$；因此达到局部最大值时也有相应不等式。
\end{proposition}
\begin{proof}
全局情形由 $T_tf(x_0)\le f(x_0)$ 后除以 $t$ 得到。局部为零时，$|T_tf(x_0)|\le\|f\|_\infty P_t(x_0,B_\varepsilon^c)=o(t)$。局部最大值时，同样将小球内的差分与小球外的 $o(t)$ 分开。若空间包含常数，保守性还给出 $1\in D(A)$ 时 $A1=0$；在非紧空间 $C_0$ 中，常数通常不属于该空间，须在核或适当扩张上理解此式。
\end{proof}

\begin{theorem}[光滑系数的尺度与速度表示]\label{thm:y13-scale-speed}
在开区间 $I=(l,r)$ 上考虑 $Lf=a(x)f''+b(x)f'$，其中 $a>0$ 连续、$b$ 连续。固定 $x_0\in I$，令
\[
 s'(x)=\exp\left(-\int_{x_0}^x\frac{b(u)}{a(u)}\,\dd u\right),\qquad
 m'(x)=\frac1{a(x)s'(x)}.
\]
则 $s$ 严格递增，且
\[
 L=D_mD_s,\qquad D_sf=\frac{f'}{s'},\qquad D_mg=\frac{g'}{m'}.
\]
调和方程 $Lf=0$ 的解为 $\alpha+\beta s$。尺度 $s$ 允许正仿射变换；相应速度测度需作逆比例调整。
\end{theorem}
\begin{proof}
由 $(\log s')'=-b/a$，直接计算
\[
 \frac1{m'}\left(\frac{f'}{s'}\right)'
 =a f''-a\frac{s''}{s'}f'=a f''+bf'.
\]
若该式为零，则 $D_sf$ 为常数，所以解恰为上述形式。将 $s$ 换为 $\gamma s+\delta$、$\gamma>0$ 时，$\dd m$ 换为 $\gamma^{-1}\dd m$，乘积算子不变。
\end{proof}

\paragraph{广义速度测度。}
令 $s$ 连续严格递增，$m$ 为局部有限且在每个开子区间有正质量的测度。广义表达式 $Lf=g$ 可用积分关系定义：$D_sf$ 有局部有界变差的适当代表，并满足
\[
 D_sf(v)-D_sf(u)=\int_{(u,v]}g\,\dd m.
\]
当 $m$ 有原子时，$D_sf$ 可以跳跃，故 $D_mD_s$ 不必是普通光滑系数的微分式。这是 Feller 广义 Laplace 表示容纳黏滞等现象的原因。原书从局部性、最大值原理及非退化条件反向构造尺度和速度；其完整反向表示定理不由本节的光滑计算替代，列入进一步阅读范围。

\paragraph{边界类型与定义域。}
选 $c\in(l,r)$，定义左端测试量
\[
 \Sigma_l=\int_{(l,c)}[s(x)-s(l)]\,m(\dd x),\qquad
 N_l=\int_{(l,c)}[s(c)-s(x)]\,m(\dd x),
\]
端点尺度取单侧极限，允许无穷值。右端相应使用 $s(r)-s(x)$ 与 $s(x)-s(c)$。Feller 边界分类定理在其一维正则扩散假设下给出
\[
 \begin{array}{c|cc}
 &N_l<\infty&N_l=\infty\\\hline
 \Sigma_l<\infty&\text{正则}&\text{流出}\\
 \Sigma_l=\infty&\text{流入}&\text{自然}
 \end{array}
\]
本节采用该分类作为外部结构定理，并计算典型例子；不声称已证明一般边界可达性的概率刻画。正则端点还须选择吸收、反射或其他允许的边界条件，才能确定闭生成元。内部的 $D_mD_s$ 表达式本身不决定半群。

\begin{example}[三种左端行为]
在 $(0,\infty)$ 上，$L=f''$ 有 $s=x$、$m'=1$，故零端两个测试量均有限，为正则端点。$L=xf''$ 有 $s=x$、$m'=1/x$，于是 $\Sigma_0<\infty$、$N_0=\infty$，零端为流出端点。对三维 Bessel 表达式 $L=f''+2f'/x$，有 $s=-1/x$、$m'=x^2$，所以 $\Sigma_0=\infty$、$N_0<\infty$，零端为流入端点。
\end{example}

\begin{exercise}[自然端点]\label{ex:y13-natural-boundary}
对 $Lf=x^2f''+xf'$、$x>0$，求尺度、速度密度与零端类型，并解释对数坐标的作用。
\end{exercise}
\begin{solution}
可取 $s=\log x$、$m'=1/x$。由于 $s(0)=-\infty$，有 $\Sigma_0=\infty$，且
\[
 N_0=\int_0^c\log(c/x)\,\frac{\dd x}{x}=\infty.
\]
所以零端为自然端点。若 $f(x)=g(\log x)$，则 $Lf=g''(\log x)$，零端对应实轴 Brownian 运动的负无穷远。
\end{solution}

\begin{exercise}[边界条件改变转移]\label{ex:y13-boundary-kernels}
令 $p_t(z)=(4\pi t)^{-1/2}e^{-z^2/(4t)}$。比较半直线上 $p_t(y-x)-p_t(y+x)$ 与 $p_t(y-x)+p_t(y+x)$ 的总质量和零端边界条件。
\end{exercise}
\begin{solution}
对 $x>0$，差核在 $y\ge0$ 上质量小于一，给出到达零端后被杀死的次 Markov 半群，边界为 Dirichlet。和核总质量为一，给出反射半群；其关于 $x$ 的零端导数为零，对应 Neumann 条件。二者内部生成式均为二阶导数，保守性与状态空间的边界处理不同。
\end{solution}

\Needspace{10\baselineskip}
\section{扩散算子的一种推广}\label{sec:yosida-13-07}
\noindent\textit{原书起始页：中译本 426；英文版 408。}

去掉 Lindeberg 条件后，短时间内发生远距离跳跃的概率可以是一阶小量。生成元因而出现积分项，通常不再具有局部性。

\begin{proposition}[有限跳跃测度的生成元]\label{prop:y13-compound-poisson}
设 $\nu$ 为 $\mathbb R^d\setminus\{0\}$ 上有限正测度，$r=\nu(\mathbb R^d)>0$、$\pi=\nu/r$。定义 $Kf(x)=\int f(x+z)\pi(\dd z)$。则 $B=r(K-I)$ 是 $C_0(\mathbb R^d)$ 上的有界生成元，
\[
 e^{tB}f=e^{-rt}\sum_{n=0}^\infty\frac{(rt)^n}{n!}K^nf,
 \qquad Bf(x)=\int[f(x+z)-f(x)]\nu(\dd z).
\]
这是保守的正收缩半群，且为强连续。
\end{proposition}
\begin{proof}
$K$ 正且范数至多一；连续性由控制收敛，无穷远消失性由先把 $\pi$ 的质量限制在大球再估计尾部得到。级数在算子范数中收敛，系数为总和等于一的非负 Poisson 权重，所以正且收缩。半群性质由指数级数得到，差商在算子范数中趋于 $B$。核可理解为以 Poisson 次数抽取独立位移；这里公式本身已构成转移核，无需先援引样本路径构造。
\end{proof}

与常系数漂移 Gaussian 半群复合，因各部分都是平移卷积而交换，得到跳跃扩散，其生成式为
\[
 Af=b\cdot\nabla f+\tfrac12Q:D^2f+
 \int[f(x+z)-f(x)]\nu(\dd z),\qquad Q\ge0.
\]
在测试函数上可直接对复合半群求导验证此式。

\begin{definition}[补偿后的跳跃项]
允许 $\nu$ 在原点附近有无限质量，但要求
\[
 \int_{\mathbb R^d\setminus\{0\}}(1\wedge|z|^2)\nu(\dd z)<\infty.
\]
对 $f\in C_c^\infty$，定义
\[
 \mathcal Jf(x)=\int\bigl[f(x+z)-f(x)-\mathbf1_{\{|z|\le1\}}z\cdot\nabla f(x)\bigr]\nu(\dd z).
\]
一阶补偿使小位移余项为 $O(|z|^2)$，所以积分绝对收敛。截断函数的改变须由漂移项的相应调整补偿，不能单独更改积分项。
\end{definition}

\paragraph{一般结构定理的边界。}
对 $\mathbb R^d$ 上具有 $C_c^\infty\subset D(A)$ 的 Feller 生成元，正最大值原理的结构表示通常写为
\begin{align*}
 Af(x)={}&-\kappa(x)f(x)+b(x)\cdot\nabla f(x)+\tfrac12Q(x):D^2f(x)\\
 &+\int\bigl[f(x+z)-f(x)-\mathbf1_{\{|z|\le1\}}z\cdot\nabla f(x)\bigr]\nu(x,\dd z),
\end{align*}
其中 $\kappa\ge0$、$Q\ge0$，每个 $x$ 的跳跃测度满足上述二阶可积条件；保守且单位函数可合法纳入定义域时杀死项为零。这是广义扩散的标准形式。原书在固定时空点的截断测度有界、紧性及差商存在条件下给出相应的子序列表示。本节将一般结构表示列为外部定理，完整构造并证明的是有限跳跃与常系数复合模型；任意给定的变系数组合是否生成 Feller 半群还需要额外存在性条件。

\begin{proposition}[最大值原理并不等于局部性]
上述跳跃积分在 $f$ 的全局最大点非正，而 $f$ 在该点邻域恒为零时积分仍可非零。对有限非零 $\nu$，相应过程不满足 Lindeberg 条件。
\end{proposition}
\begin{proof}
最大点处梯度为零，每个差分非正。若 $f$ 远处为正，且跳跃测度能跳到该处，则 $\mathcal Jf(x)>0$，显示非局部性。选 $\varepsilon>0$ 使 $\nu\{|z|>\varepsilon\}>0$，Poisson 级数的一次跳跃项给出
\[
 \liminf_{t\downarrow0}t^{-1}P_t(x,\{|y-x|>\varepsilon\})
 \ge\nu\{|z|>\varepsilon\}>0.
\]
故局部扩散所需的小时间条件失败。
\end{proof}

\begin{exercise}[对称跳跃的符号]\label{ex:y13-jump-symbol}
在实线上取 $\nu=\tfrac r2(\delta_a+\delta_{-a})$、$a\ne0$。写出 $A$，并求它在有界 Fourier 模 $e^{i\xi x}$ 上的形式作用。
\end{exercise}
\begin{solution}
$Af(x)=\tfrac r2[f(x+a)+f(x-a)-2f(x)]$。形式代入得到 $Ae^{i\xi x}=r(\cos(a\xi)-1)e^{i\xi x}$。这些模不在 $C_0$ 中，等式在有界函数扩张或符号计算意义下理解；乘以测试截断后仍可用于辨认跳跃项，不能当作 $C_0$ 的本征向量结论。
\end{solution}

\begin{exercise}[小跳跃的补偿]\label{ex:y13-small-jumps}
实线上 $\nu(\dd z)=\mathbf1_{0<|z|<1}|z|^{-1-\alpha}\,\dd z$，$0<\alpha<2$。验证二阶可积条件，并说明何时未补偿的一阶绝对估计足够。
\end{exercise}
\begin{solution}
$\int z^2\,\dd\nu=2\int_0^1r^{1-\alpha}\,\dd r=2/(2-\alpha)<\infty$，所以补偿后的 Taylor 余项可积。一阶绝对估计要求 $\int|z|\,\dd\nu=2\int_0^1r^{-\alpha}\,\dd r<\infty$，恰在 $\alpha<1$ 时成立。对称性可允许主值抵消，但主值与绝对收敛不能混称。
\end{solution}

\Needspace{10\baselineskip}
\section{Markov 过程和位势}\label{sec:yosida-13-08}
\noindent\textit{原书起始页：中译本 429；英文版 410。}

设 $T_t$ 为 Banach 空间上的有界强连续半群，生成元 $A$，此节先假定为正收缩。记
\[
 R_\lambda=(\lambda I-A)^{-1}
 =\int_0^\infty e^{-\lambda t}T_t\,\dd t,\qquad\lambda>0.
\]
右侧对每个向量是 Bochner 积分。核表示下它也是折扣占用算子。$h\in D(A)$ 满足 $Ah=0$ 称为调和元素；等价地 $T_th=h$，或 $\lambda R_\lambda h=h$ 对所有 $\lambda>0$ 成立。

\Needspace{12\baselineskip}
\begin{proposition}[位势与 Poisson 方程]\label{prop:y13-poisson}
若 $R_\lambda g$ 在 $\lambda\downarrow0$ 时有范数极限 $v$，则 $v\in D(A)$ 且 $Av=-g$。特别地，若 $\int_0^\infty\|T_tg\|\,\dd t<\infty$，则
\[
 v=\int_0^\infty T_tg\,\dd t=\lim_{\lambda\downarrow0}R_\lambda g.
\]
反向，$Av=-g$ 的解可以相差调和元素，不能仅凭方程断言唯一。
\end{proposition}
\begin{proof}
$AR_\lambda g=\lambda R_\lambda g-g\to-g$，而 $R_\lambda g\to v$，生成元闭性给出结论。绝对可积时，Bochner 控制收敛给出积分极限。两个解之差属于 $\ker A$，故最后的限制必要。
\end{proof}

\begin{theorem}[正格中的 Riesz 型分解]\label{thm:y13-riesz-potential}
设空间为 Dedekind $\sigma$ 完备 Banach 格，范数序列序连续。若 $u\in D(A)$、$u\ge0$、$Au\le0$，则
\[
 h=\lim_{\lambda\downarrow0}\lambda R_\lambda u,\qquad
 p=u-h=\lim_{\lambda\downarrow0}R_\lambda(-Au)
\]
均以范数收敛，$h\ge0$ 调和，$p\ge0$ 为位势。$h$ 是所有不超过 $u$ 的调和元素中的最大者。相应地，若次调和 $x\in D(A)$、$Ax\ge0$ 有调和上界，则它有最小调和上界，并分解为该上界加一个非正位势。
\end{theorem}
\begin{proof}
恒等式 $\lambda R_\lambda u-u=R_\lambda Au\le0$ 给出 $0\le\lambda R_\lambda u\le u$。对 $\lambda>\mu>0$，预解式恒等式给出
\[
 (\lambda R_\lambda-\mu R_\mu)u
 =-(\lambda-\mu)R_\lambda R_\mu Au\ge0.
\]
故 $\lambda\downarrow0$ 时该族递减。先沿 $\lambda_n\downarrow0$ 取格下确界，序连续范数使序列趋于 $h$；任意中间参数由单调夹逼得到同一范数极限。由于
\[
 A(\lambda R_\lambda u)=\lambda(\lambda R_\lambda u-u)\longrightarrow0,
\]
闭性给出 $Ah=0$。再用 $R_\lambda(-Au)=u-\lambda R_\lambda u$ 得 $p$ 及其位势表示。若 $k\le u$ 调和，则 $k=\lambda R_\lambda k\le\lambda R_\lambda u$，取极限得 $k\le h$。次调和情形对 $u=H-x\ge0$ 应用以上结论，其中 $H$ 是任一调和上界，即得最小调和上界。
\end{proof}

此定理是原书次调和分解的一个明确 Banach 格版本；一般局部凸序空间的弱单调完备假设不在此处另作推广。特别地，$C_0(S)$ 的一致范数通常不满足序连续假设，不能把此定理未经检查就用于所有连续函数位势。

\begin{example}[欧氏热核的暂留与再返界限]\label{exm:y13-green}
取生成元 $\Delta$。当 $d\ge3$ 且 $g\in C_c(\mathbb R^d)$，有 $\int_0^\infty\|T_tg\|_\infty\,\dd t<\infty$：小时间用收缩性，大时间用 $(4\pi t)^{-d/2}\|g\|_1$。因而
\[
 Vg(x)=\int_{\mathbb R^d}G_d(x-y)g(y)\,\dd y,\qquad
 G_d(z)=\frac{|z|^{2-d}}{(d-2)|\mathbb S^{d-1}|}.
\]
该核由 $\int_0^\infty(4\pi t)^{-d/2}e^{-|z|^2/(4t)}\,\dd t$ 换元计算，$-\Delta Vg=g$ 按生成元及分布意义成立；不据此声称一般连续 $g$ 的位势有连续二阶导数。对 $d=1,2$ 和非零非负 $g\in C_c$，固定 $x$ 时大时间的 $T_tg(x)$ 与 $t^{-d/2}\int g$ 同阶，故占用积分发散。该例说明无折扣位势可能不存在，虽每个 $R_\lambda$ 都存在。
\end{example}

\begin{exercise}[平衡态的障碍]\label{ex:y13-potential-obstruction}
在保概率 Markov 半群上，证明 $g=1$ 没有有限范数 Abel 位势；并证明 $Av=-g$ 有可积定义域中的解时必有 $\int g\,\dd m=0$。
\end{exercise}
\begin{solution}
$R_\lambda1=1/\lambda$，范数随 $\lambda\downarrow0$ 发散。若 $v\in D(A)\subset L^1$，保积分性对差商取极限给出 $\int Av=0$，所以 $Av=-g$ 要求 $\int g=0$。非零均值是 Poisson 方程的平衡态障碍。
\end{solution}

\begin{exercise}[有限链上的调和与位势]\label{ex:y13-two-state-potential}
令 $A=\left(\begin{smallmatrix}-a&a\\b&-b\end{smallmatrix}\right)$、$a,b>0$。求不变概率、平均投影，并在零均值子空间求位势。
\end{exercise}
\begin{solution}
不变概率为 $\pi=(b,a)/(a+b)$，$Qf=(\pi f)1$。零均值子空间由 $(a,-b)$ 张成，$A$ 在其上为 $-(a+b)I$，故位势 $Vg=g/(a+b)$。全空间上常数调和分量阻止无折扣逆，但投到零均值空间后可逆。
\end{solution}

\Needspace{10\baselineskip}
\section{抽象位势算子和半群}\label{sec:yosida-13-09}
\noindent\textit{原书起始页：中译本 430；英文版 411。}

本节不要求有次序或正性。$X$ 为实或复 Banach 空间，$T_t$ 为收缩强连续半群，$A$ 为其生成元。把连续对偶伴随写为 $V^*$，配对按对偶线性泛函理解。

\noindent\textit{原书校注：}末尾与 Hunt 理论比较的段落，把预解式位势的极限下标排成趋于无穷，与本节式 (7) 的零端极限不符。本稿统一采用 $\lambda\downarrow0$；实际上收缩半群满足 $\|R_\lambda\|\le1/\lambda$，若令 $\lambda\to\infty$，未乘 $\lambda$ 的预解式只能趋于零。

\begin{theorem}[零端 Abel 极限与抽象位势]\label{thm:y13-abstract-potential}
下列条件等价：$A$ 单射且 $\overline{R(A)}=X$；对每个 $x\in X$ 有 $\lambda R_\lambda x\to0$（$\lambda\downarrow0$）。此时
\[
 V=-A^{-1},\quad D(V)=R(A),\quad R(V)=D(A)
\]
为闭、定义域及值域均稠密的算子，称为抽象位势，满足
\[
 D(V)=\{x:\lim_{\lambda\downarrow0}R_\lambda x\text{ 在范数中存在}\},\qquad
 Vx=\lim_{\lambda\downarrow0}R_\lambda x.
\]
\end{theorem}
\begin{proof}
若 $x=Ay$，则 $\lambda R_\lambda x=\lambda(\lambda R_\lambda y-y)\to0$。因 $\|\lambda R_\lambda\|\le1$，值域稠密时此极限扩展到全部 $X$。反向若全部极限为零，则 $Ah=0$ 蕴含 $\lambda R_\lambda h=h=0$；又 $x-\lambda R_\lambda x=-AR_\lambda x\in R(A)$ 趋于 $x$，故值域稠密。

取 $x=-Ay$，则 $R_\lambda x=y-\lambda R_\lambda y\to y=Vx$。反之，若 $R_\lambda x\to y$，闭性与 $AR_\lambda x=\lambda R_\lambda x-x\to-x$ 给出 $Ay=-x$，所以 $x\in R(A)$。逆算子的图由 $A$ 的图交换坐标得到，故 $V$ 闭；两个稠密性由假设及生成元定义域稠密得到。
\end{proof}

\begin{theorem}[位势的双侧范数刻画]\label{thm:y13-potential-characterization}
设 $V$ 为闭线性算子，$D(V)$ 与 $R(V)$ 均在 $X$ 中稠密。则 $V$ 是某收缩强连续半群的抽象位势，当且仅当对每个 $\lambda>0$ 有
\begin{align}
 \|\lambda Vx\|&\le\|x+\lambda Vx\|&& (x\in D(V)),\label{eq:y13-potential-primal}\\
 \|\lambda V^*\phi\|&\le\|\phi+\lambda V^*\phi\|&& (\phi\in D(V^*)).\label{eq:y13-potential-dual}
\end{align}
此时 $A=-V^{-1}$，且 $R_\lambda=V(I+\lambda V)^{-1}$。
\end{theorem}
\begin{proof}
先证明充分性。令 $B_\lambda=I+\lambda V$。第一不等式还给出 $\|x\|\le2\|B_\lambda x\|$，所以 $B_\lambda$ 单射。若 $B_\lambda x_n$ 收敛，同一估计使 $x_n,Vx_n$ 均 Cauchy，闭性使极限仍在图中。因此 $R(B_\lambda)$ 闭。其零化子为 $\ker(I+\lambda V^*)$；第二不等式使该核为零。Hahn--Banach 于是给出值域稠密，结合闭性得满射。

令 $J_\lambda=VB_\lambda^{-1}$，则 $\|\lambda J_\lambda\|\le1$。直接在 $D(V)$ 上代入并利用 $V$ 与其有界预解算子的交换关系可得
\[
 J_\lambda-J_\mu=(\mu-\lambda)J_\lambda J_\mu.
\]
对 $y=Vx$，$\lambda J_\lambda y=y-J_\lambda x\to y$（$\lambda\to\infty$），因 $R(V)$ 稠密，这个极限对全部 $y$ 成立。伪预解式的核与参数无关，故其核为零；又 $R(J_\lambda)=R(V)$ 稠密。因此第八章的伪预解式构造给出稠密定义的闭算子 $A$，其预解式为 $J_\lambda$。$\|\lambda J_\lambda\|\le1$ 蕴含各幂的 Hille--Yosida 估计，故 $A$ 生成收缩强连续半群。

由 $J_\lambda(B_\lambda x)=Vx$ 得 $AVx=-x$。$V$ 单射也可从 $\ker J_\lambda=0$ 得到，故 $A=-V^{-1}$、$R(A)=D(V)$ 稠密。上一条定理将 $V$ 识别为该半群的抽象位势。

必要性方面，若 $V=-A^{-1}$，预解式给出
\[
 B_\lambda^{-1}=I-\lambda R_\lambda,\qquad
 VB_\lambda^{-1}=R_\lambda.
\]
所以 $\lambda VB_\lambda^{-1}$ 为收缩，得到第一不等式。闭稠密算子的伴随逆关系给出 $B_\lambda^{*-1}=(B_\lambda^{-1})^*$，且 $\lambda V^*B_\lambda^{*-1}=(\lambda VB_\lambda^{-1})^*$ 仍为收缩，得到第二不等式。这里未假定伴随半群在整个 $X^*$ 上强连续。
\end{proof}

\begin{corollary}[Hilbert 空间的增生表述]
在 Hilbert 空间中，上述两条范数不等式分别等价于
\[
 \operatorname{Re}\langle Vx,x\rangle\ge0,\qquad
 \operatorname{Re}\langle V^*y,y\rangle\ge0.
\]
因此闭、定义域和值域稠密的 $V$ 是抽象位势，当且仅当 $V,V^*$ 均增生。
\end{corollary}
\begin{proof}
展开 $\|x+\lambda Vx\|^2-\|\lambda Vx\|^2=\|x\|^2+2\lambda\operatorname{Re}\langle Vx,x\rangle$。它对所有 $\lambda>0$ 非负恰好等价于增生性；伴随同理。
\end{proof}

\begin{example}[乘法位势的定义域]
在 $L^2(0,1)$ 上取 $(T_tf)(s)=e^{-ts}f(s)$。则 $Af=-sf$，它单射但逆不有界，
\[
 (Vg)(s)=g(s)/s,\qquad D(V)=\{g:g(s)/s\in L^2(0,1)\}.
\]
$R_\lambda g=g/(\lambda+s)$；$\lambda R_\lambda g\to0$ 由控制收敛成立。只有满足上述定义域条件的 $g$ 才有位势范数极限，不能把所有有界预解式的零端极限称为处处定义的有界算子。
\end{example}

\begin{exercise}[Abel 位势不要求普通积分收敛]\label{ex:y13-abel-oscillation}
在 $X=\mathbb C$ 上令 $T_tz=e^{it}z$。求 $A,V$，并比较 $\int_0^R T_tz\,\dd t$ 与 $\int_0^\infty e^{-\lambda t}T_tz\,\dd t$ 的极限。
\end{exercise}
\begin{solution}
$A=iI$、$V=-1/i=iI$。普通积分为 $(e^{iR}-1)z/i$，对非零 $z$ 随 $R\to\infty$ 不收敛；折扣积分为 $z/(\lambda-i)$，随 $\lambda\downarrow0$ 趋于 $iz$。故抽象位势容纳 Abel 可和但普通占用积分不收敛的情形，也不要求格意义下的正性。
\end{solution}

\begin{exercise}[乘法位势的临界指数]\label{ex:y13-potential-domain}
在上例取 $g(s)=s^\beta$、$\beta>-1/2$。判断何时 $g\in D(V)$，并说明不在定义域时预解式范数的行为。
\end{exercise}
\begin{solution}
$g\in D(V)$ 等价于 $\int_0^1s^{2\beta-2}\,\dd s<\infty$，即 $\beta>1/2$。若 $\beta\le1/2$，则 $\|R_\lambda g\|_2^2=\int_0^1s^{2\beta}/(\lambda+s)^2\,\dd s\uparrow\infty$，不可能有范数极限，虽对每个正 $\lambda$ 都是良定义的预解式值。
\end{solution}


\Needspace{6\baselineskip}
\section*{本章参考说明与阅读衔接}
直接来源为吉田中英文第六版。个别遍历证明以有限极大不等式重组；不变测度分解在紧可度量 Feller 情形展开；扩散部分以局部矩、尺度速度与显式跳跃模型解释生成元。一般比值遍历、非紧空间的完整逐点分解、广义尺度速度的反向刻画、一般 Feller 边界分类及变系数跳跃表示的证明边界均明确记录，不把模型计算等同于全部结构定理。下一章将转向发展方程的积分与初值问题。
